\documentclass{article}
\usepackage[utf8]{inputenc}
\usepackage[T1]{fontenc}
\usepackage{geometry}
\geometry{
    a4paper,
    left=2.5cm,
    right=2.5cm,
    top=2.5cm,
    bottom=2.5cm
}
\usepackage{graphicx}
\usepackage{float}
\usepackage{hyperref}
\usepackage[english]{babel}
\usepackage{amsmath}
\usepackage{amssymb}
\usepackage{array}
\usepackage{booktabs}
\usepackage{caption}
\usepackage{subcaption}
\usepackage{enumitem}
\usepackage{fancyhdr}
\usepackage{titlesec}
\usepackage{url}
\usepackage{textcomp}
\usepackage{xcolor}
\usepackage{listings}
\usepackage{multirow}
\usepackage{longtable}
\usepackage{tabularx}
\usepackage{times}

% Set up page style
\pagestyle{fancy}
\fancyhf{}
\rfoot{\thepage}
\renewcommand{\headrulewidth}{0pt}

% Title formatting
\titleformat{\chapter}[display]{\normalfont\bfseries}{\chaptertitle\thechapter}{0pt}{\Huge}
\titlespacing*{\chapter}{0pt}{50pt}{40pt}

% Document info
\title{Template Document}
\author{Academic Support Agent}
\date{\today}

\begin{document}
\maketitle

\title{Sample Article for\break {Oxford University Press~Journals}}

\author{%
\name{First Author}
\address{Institute or Organization, Department, City, State,\\
Zip Code, Country}
\email{e-mail address}
\and
\name{Second Author}
\address{Institute or Organization, Department, City, State,\\
Zip Code, Country}
\email{e-mail address}}

\abstract{This sample is a guideline for preparing technical papers using \LaTeX\
for manuscript submission. It contains the documentation for your \LaTeX\
Class file, which implements the layout for your manuscript for all Journals of OUP.
This sample article uses a class file named \texttt{ouparticle.cls} that all authors
need to use for their manuscript preparation. It is similar in use to the \texttt{article}
class file of \LaTeX, but has some extra fields in the preamble and some extended commands for
other parts of the article.}

\date{\today}

\keywords{word1; word2; word3; and word4}

\maketitle


\section{Introduction}
\label{sec1}

It is assumed that the author is familiar with either plain
\TeX, \AmS-\TeX{} or a standard \LaTeX\ setup and, hence,
only the essential points are described in this document.
Nevertheless, we hope that this document is generally sufficient
for describing the requirements for preparation of
manuscripts. For more details, please see the \textit{\LaTeX{} User's Guide} or
\textit{The not so short introduction to \LaTeXe}.


\section{Installation}
\label{sec2}

Provided with \verb+ouparticle.cls+ are the files
\verb+sample.tex+ (this document explains the various
features of \verb+ouparticle.cls+) and \verb+sample.pdf+
(how the output using \verb+sample.tex+ should be). Your
paper can be compiled with standard \LaTeX,
preferably with the current \LaTeXe\ version. It will probably work
with older versions of \LaTeXe; however, this has not been
tested. The file \verb+ouparticle.cls+ needs to be copied
into a directory where \TeX\ looks for input files. The other files
need to be kept as a reference while preparing
your manuscript. Please use the predefined commands from \verb+sample.tex+ for title,
authors, abstract, body, etc.


\section{Preparing your manuscript}
\label{sec3}

\subsection{General guidelines}
\label{sec3.1}

\begin{enumerate}
\item
\LaTeX\ and \AmS-\LaTeX\ provide a rich set of commands for all common, important
features of your paper. Use them; avoid definitions and use of custom commands.

\item
There is no need to redefine any \TeX, \LaTeX\ or \AmS-\LaTeX\ commands.

\item
Avoid direct formatting for headings cleanly set as section headings.

\item
Use \LaTeX\ commands for font changes. For example: use \verb+\textbf{phrase}+,
not \verb+{\bf phrase}+; use \verb+\mathcal{C}+, not \verb+{\cal C}+; etc.
\end{enumerate}


\subsection{How to start with \texttt{ouparticle.cls}}
\label{sec3.2}

Before you type anything that actually appears in the paper, you need to
include a \verb+\documentclass{ouparticle}+ command at the very beginning,
and then the two commands that have to be part of any \LaTeX\ document,
\verb+\begin{document}+ at the start and \verb+
\end{document}